\chapter{Verification and Evaluation Design}
\label{ch:verification}
\section{Reference Hierarchy}
The evaluation proceeds from exact checks to independently discretized
comparisons. This hierarchy separates equation implementation, numerical
approximation, and physical applicability \cite{roache1998}.

For plane Poiseuille flow in a straight channel of half-width $h$,
\begin{equation}
 u=1-(y/h)^2,\qquad v=0,\qquad p=-2x/h^2+C.
 \label{eq:poiseuille}
\end{equation}
This profile checks both FEM and a separately trained straight-channel PINN.
Velocity profiles are compared on a transverse section; relative pressure
errors use points spanning the channel length, because exact pressure has
no transverse variation on a single section.

In a uniform flow, the Stokes relaxation solution provides an independent
particle-ODE check,
\begin{equation}
 \vect v_p(t)=\vect u+
 [\vect v_p(0)-\vect u]\exp(-t/\St).
 \label{eq:relaxation}
\end{equation}
Algebraic drag tests quantify whether ForceNet reproduces its specified law;
they do not validate that law against real particles near a wall.

\section{Field and Trajectory Metrics}
The reported relative field error is a discrete sample norm,
\begin{equation}
 E_u=\frac{\sqrt{\sum_i\lVert\widehat{\vect u}_i-\vect u_i^{\rm ref}\rVert^2}}
 {\sqrt{\sum_i\lVert\vect u_i^{\rm ref}\rVert^2}}.
 \label{eq:relative-error}
\end{equation}
Pressure uses the same construction after each pressure sample set is
mean-centered to remove gauge differences. Interior sampling is in channel
coordinates; the resulting norm is not an exact area-weighted integral.
Maximum velocity error, boundary RMS error, PDE residuals and section flux
provide complementary information. FEM's integrated divergence norm and
PINN's sampled continuity RMS have different definitions and should not be
compared as if they were identical measures.

For a pair of paths, the position RMS is calculated at 150 common-time
samples up to $T_c=\min(T_{\rm PINN},T_{\rm ref})$. Its ratio to width $W$
is reported alongside $T_c/\max(T_{\rm PINN},T_{\rm ref})$ and endpoint-class
agreement. Restricting the time interval can hide disagreement after one
path ends; the coverage and class measures prevent that truncation from
being interpreted as full-path agreement.

\begin{table}[!htbp]
\centering\small
\caption{Numerical screening criteria used to organize the proposal.}
\label{tab:criteria}
\begin{tabularx}{\textwidth}{@{}p{.43\textwidth}Y@{}}
\toprule Quantity & Screening requirement\\\midrule
Channel velocity and pressure & Relative sampled error below 5\% each\\
Straight-channel PINN & Velocity and pressure error below 5\% each\\
Boundary velocity & Dimensionless RMS below 0.02\\
Section mass flux & Maximum relative deviation below 5\%\\
Reference mesh difference & Penultimate-to-finest velocity difference below 2\%\\
Reference trajectory interpolator & Velocity discrepancy below 2\%\\
Force surrogate & Held-out normalized relative error below 5\%\\
Baseline trajectory position & RMS below 5\% of width for every compared seed\\
Baseline endpoint classes & Agreement for every compared seed\\
Common trajectory time & More than 95\% for every compared seed\\
Separate ABC benchmark & Velocity and pressure errors below 5\% each\\
\bottomrule
\end{tabularx}
\end{table}
These are numerical screening choices, not universal engineering tolerances.
The current records are a developmental baseline; this proposal does not
claim that the listed thresholds were preregistered before all preliminary
inspection. Final acceptance criteria and held-out cases will be frozen
before the proposed comparison experiments begin.

\section{Separate Three-Dimensional Benchmark}
The exact ABC field, with unit amplitudes, is
\begin{align}
 \vect u&=(\sin z+\cos y,\ \sin x+\cos z,\ \sin y+\cos x),\nonumber\\
 p&=-\Rey\lVert\vect u\rVert^2/2,\qquad \vect f=\vect u,
 \label{eq:abc}
\end{align}
on $[0,2\pi]^3$. It satisfies the forced steady momentum equation
$\Rey(\vect u\cdot\nabla)\vect u=-\nabla p+\nabla^2\vect u+\vect f$.
Indeed, $\nabla\cdot\vect u=0$, $\nabla\times\vect u=\vect u$, and
$\nabla^2\vect u=-\vect u$. The convective term is the gradient of
$\lVert\vect u\rVert^2/2$, which is canceled by the pressure gradient.

A separate 3D network predicts all three velocity components and pressure;
known velocity is prescribed on the cube faces, and one pressure value fixes
the gauge. The exact interior field is reserved for evaluation. Particle-free
tracer paths use the periodically extended exact or learned velocity. This
tests genuine 3D derivatives and trajectories, but not a finite-depth curved
channel with sidewalls. The extruded planar channel illustration has an
entirely different role and is identified as visualization only.

\section{Uncertainty and Physical Validation}
The current error budget includes reference mesh and geometry error,
interpolator error, PINN approximation, optimizer variability and ODE error.
Two seeds and a three-mesh comparison do not constitute a formal uncertainty
analysis or a grid-convergence index. Temporal and event-location refinement
are required in the proposed work. At present, no measurement dataset has
been supplied; physical validation and uncertain particle properties remain
open requirements before a device-level claim is justified.
